\section{Primitive inputs to the record and report bounds}
\label{app:primitive-inputs}

The composition theorems use operational distinguishing bounds: a record
commitment comparison and a single immutable report row. This appendix fixes
the concrete parameter conventions. The public hash and compression oracles
are idealized in the stated games; the source construction uses the specified
SHA3 and HMAC-SHA256 algorithms.

\subsection{Hidden-input random-oracle replacement}

For a hidden-input entropy parameter $n$, define
\[
 \Phi_n(d,Q)=\min\!\left\{1,
 4\sqrt{2d^2Q\,2^{-n}}+4\sqrt{2Q^2d\,2^{-n}}\right\}.
\]
Don, Fehr and Huang's Theorem~2~\cite{dfh2022} supplies this bound when
private oracle inputs have point probability at most $2^{-n}$ and distinct
effective private queries encode distinct inputs for every secret value.
Exact repeats are cached before that interface is queried. The two budgets
count private classical queries and all ordinary public-oracle queries;
public queries may be quantum.

For image $i$, let $K_i^{\rm hdr}$ count original intent prefixes and let
$d_i\le1+5K_i^{\rm hdr}$ count its distinct private digest inputs. There is
one image-head domain and five intent-digest domains. With an independent
192-bit seal nonce, the fixed-input encoding and per-nonce confinement
conditions give the retained record bound
\[
 \epsilon_{\rm com}
 \le\min\left\{1,\sum_{i=1}^{M}\Phi_{192}(d_i,Q_i)\right\}.
\]
Here $d_i=0$ for known, atom-free rows. The budget $Q_i$ includes the
entire derived distinguisher, all other real rows and the final test.
The nonce remains private before the declared acquisition stop. Local
image encryption and intent authentication introduce no additional shared
SHA3 calls. This comparison is made after private-wire equality has been
replaced by structural equality; its separate collision term is the one in
Corollary~\ref{cor:real-to-opaque}. The exact encoding domains and fixed-source
arguments are identified in Appendix~\ref{app:opaque-artifacts}.

\subsection{One immutable HMAC row}

For the row experiment in Section~\ref{sec:report-composition}, let
\[
 h:\{0,1\}^{512}\times\{0,1\}^{256}\longrightarrow\{0,1\}^{256}
\]
be the public compression oracle, with arguments in message-block,
chaining-state order. A 32-byte key $K$ determines the two pad blocks
\[
 I(K)=(K\Vert0^{32})\oplus\mathtt{36}^{64},\qquad
 O(K)=(K\Vert0^{32})\oplus\mathtt{5c}^{64}.
\]
Zero exponents in these encodings count bytes. The ranges are disjoint and
each encoding is injective in $K$. Cache the two evaluations at the fixed
SHA-256 initial state. All oracle-dependent preparation is executed and
counted; no oracle-correlated advice is supplied for free.

Let $L$ be the maximum message length in bytes and put
\[
 L\le2^{61}-65,\qquad
 \ell=1+\left\lceil\frac{L+9}{64}\right\rceil.
\]
The SHA-256 padding encodes $8(64+|m|)$ bits, retaining the already consumed
64-byte key-pad block. Its suffix is determined by $|m|$ and it is injective
within the stated length domain. The parameter $\ell$ bounds the remaining
row compression work after the two key-pad calls.

Write $E_{\rm NMAC}(T,q,\ell)$ for a valid NMAC distinguishing bound with two
independent 256-bit inner/outer keys, the prescribed injective padding, at
most $T$ keyed calls and $q$ public compression calls. Hosoyamada and Iwata's
Proposition~6~\cite{hmacqrom2021} gives the asymptotic form
\[
 E_{\rm NMAC}(T,q,\ell)
 =O\!\left(\sqrt{\frac{(q+T)^3\ell^5}{2^{256}}}\right)
\]
under its stated padding and query conditions. The implicit constant remains
part of that published bound. Our message-length convention retains HMAC's
consumed-block offset.

The two key-pad replacements use $\Phi_{256}(2,q+\ell T)$. The subsequent
NMAC replacement gives an admissible row bound
\[
 \epsilon_{\rm row}(T,q,L)
 =\min\!\left\{1,
 \Phi_{256}(2,q+\ell T)+E_{\rm NMAC}(T,q,\ell)\right\}.
\]
This is a restricted single-transcript application of the primitive results.
The ideal NMAC function is queried on one original transcript only, so its
first reply is a uniform tag and exact repeats retrieve the same value.
Report nonzero-ID checks act on that reply in both games. All other real
HMAC work remains in $q$; the multi-report theorem supplies the dependency
order and its actual disclosure residual.
